% ==========================================
% Document Content for Radiation Safety Manual
% Translation of Spanish Text to LaTex
% Spelling checked and corrected
% ==========================================

\section*{PRÁCTICA No. 8}

\textbf{MANEJO DE UN EQUIPO DE GAMMAGRAFÍA}

\textbf{OBJETIVO:} Una vez terminada satisfactoriamente esta práctica, los alumnos deben estar en posibilidad absoluta de armar, utilizar, desarmar y dar mantenimiento adecuado a un equipo completo de gammagrafía.

\textbf{MATERIAL:}
\begin{itemize}
    \item[1] Contenedor de Gammagrafía. (Con todos sus Accesorios).
    \item[1] Fuente simulada.
    \item[1] Equipo portátil Detector de Radiación tipo Geiger-Müller.
    \item[1] Juego completo de Croquis relacionados con el equipo de que se trate, así como de todos sus componentes.
\end{itemize}

\textbf{DESARROLLO:} Se le entrega al alumno el Equipo Completo de Gammagrafía (incluyendo un Equipo Geiger-Müller). El alumno, por sí solo, efectuará los diferentes pasos involucrados en el manejo de un equipo de esta naturaleza, desde el propio armado inicial, pasando por el proceso habitual de radiografiado, hasta su desmontaje, posterior mantenimiento final de limpieza y almacenamiento seguro del mismo.

Se analizarán las fallas y se encontrará la solución adecuada. Al finalizar con la última etapa de la práctica, el instructor hará una evaluación de los conocimientos y la eficiencia que en el manejo del equipo, mostró el alumno.

Esta acción, en el caso de que haya varias personas participando en la práctica, deberá repetirse por cada alumno presente.

% End of Document Content for Radiation Safety Manual
% ==========================================